%% LyX 2.4.4 created this file.  For more info, see https://www.lyx.org/.
%% Do not edit unless you really know what you are doing.
\documentclass[english]{article}
\usepackage[T1]{fontenc}
\usepackage[latin9]{inputenc}
\usepackage{enumitem}
\usepackage{amsmath}
\usepackage{amsthm}
\usepackage{geometry}
\geometry{verbose,tmargin=1in,bmargin=1in,lmargin=1in,rmargin=1in}

\makeatletter
%%%%%%%%%%%%%%%%%%%%%%%%%%%%%% Textclass specific LaTeX commands.
\numberwithin{equation}{section}
\numberwithin{figure}{section}
\newlength{\lyxlabelwidth}      % auxiliary length 

%%%%%%%%%%%%%%%%%%%%%%%%%%%%%% User specified LaTeX commands.
\usepackage{fancyhdr}
\usepackage{lastpage}
\pagestyle{fancy}
\usepackage{enumitem}
\usepackage{ifthen}
\everymath=\expandafter{\the\everymath\displaystyle}
%\DeclareMathSizes{10}{10}{10}{10}
\usepackage{multicol}

\usepackage{tikz}
\usepackage{tikz-3dplot}
\usetikzlibrary{calc,arrows}

\usetikzlibrary{patterns}
\usetikzlibrary{plotmarks}
\usepackage{pgfplots}
\pgfplotsset{compat=newest}
\usepgfplotslibrary{polar}

\usepackage{calc}
\usepackage[nomessages]{fp}

\usepackage{datetime}
% Added by lyx2lyx
\setlength{\parskip}{\smallskipamount}
\setlength{\parindent}{0pt}

\makeatother

\usepackage{babel}
\begin{document}
\renewcommand{\author}{Dr.\,James Ashe}

\newcommand{\classNum}{331}

\newcommand{\classSec}{A}

\newcommand{\examType}{Exam 4 review}

\renewcommand{\date}{\formatdate{20}{11}{2013}}

%%First page's header%%%%%%%%%%%%%%%%%%%%%%
\fancypagestyle{firststyle} %%header settings for first pagr
{
	\fancyhead[L]{MTH \classNum{}(\classSec{}) \author{}\\
	\textbf{\examType{}} ~\date{}}
	\fancyhead[C]{\textbf{Name:}}
	%\fancyhead[R]{\textbf{Name:\hspace{5cm}}}
	\setlength{\headsep}{14pt}
}
%%%%%%%%%%%%%%%%%%%%%%%%%%%%%%%%%%%%%%%%%%

%%Next page's header%%%%%%%%%%%%%%%%%%%%%%%
\setlist{nolistsep}
\lhead{MTH \classNum{}(\classSec{})} 
\chead{\textbf{Name:}} 
\cfoot{\thepage{}~of~\pageref{LastPage}} 
\setlength{\headsep}{5pt} 
%%%%%%%%%%%%%%%%%%%%%%%%%%%%%%%%%%%%%%%%%%%

%%Various settings; see the preamble for other settings%%%%%
%\setenumerate[0]{leftmargin=12pt,itemindent=0pt,labelwidth=0pt}
%\setlist{topsep=.5pt}
\renewcommand{\labelenumi}{\textbf{\arabic{enumi}.}}
\renewcommand{\labelenumii}{\textbf{\alph{enumii}.}}
\thispagestyle{firststyle}
%%%%%%%%%%%%%%%%%%%%%%%%%%%%%%%%%%%%%%%%%%

\newcommand{\xmax}{0}
\newcommand{\xmin}{-\xmax}
\newcommand{\xstep}{0}
\newcommand{\ymax}{0}
\newcommand{\ymin}{-\ymax}
\newcommand{\ystep}{0} 
\newcommand{\dspace}{\vspace{1cm}}
\newcommand{\xa}{0}
\newcommand{\xb}{0}
\newcommand{\ya}{1}
\newcommand{\yb}{1}

\newcommand{\xspace}{\vspace{1cm}}

\subsubsection*{13.1 Planes and Surfaces}
\begin{itemize}
\item Find equations of the plane, given a normal vector and a point; given
a set of points.

\begin{itemize}
\item [\textbf{ex.}]$P_{0}(1,0,-3)$; $\mathbf{n}=\left\langle 1,-1,2\right\rangle $
\item [\textbf{ex.}]$(2,-1,4)$, $(1,1,-1)$, and $(-1,3,1)$
\end{itemize}
\end{itemize}

\subsubsection*{13.2 Planes and Surfaces}
\begin{itemize}
\item Find the domain and range of a given function.

\begin{itemize}
\item [\textbf{ex.}]$f(x,y)=\sin\left(\frac{x}{y}\right)$, $f(x,y)-\sqrt{x-2y+4}$
\end{itemize}
\item Be able to read level curves.
\end{itemize}

\subsubsection*{13.3 Limits and Continuity}
\begin{itemize}
\item Limits of two-variable continuous functions.

\begin{itemize}
\item [\textbf{ex.}]\textbf{ ${\displaystyle \lim_{(x,y)\rightarrow(2,0)}\left(xy^{5}-x\right)}$},
\textbf{${\displaystyle \lim_{(x,y)\rightarrow(e^{2},4)}\ln\sqrt{xy}}$}
\end{itemize}
\item Limits at the boundary points of two-variable functions. 

\begin{itemize}
\item [\textbf{ex.}]\textbf{${\displaystyle \lim_{(x,y)\rightarrow(6,2)}\frac{x^{2}-3xy}{y-3y}}$},
\textbf{${\displaystyle \lim_{(x,y)\rightarrow(1,2)}\frac{\sqrt{y}-\sqrt{x+1}}{y-x-1}}$}
\end{itemize}
\item Nonexistence of limits using the Two-Path Test.

\begin{itemize}
\item [\textbf{ex.}]\textbf{${\displaystyle \lim_{(x,y)\rightarrow\left(0,0\right)}\frac{x+2y}{x-2y}}$},
use lines $x=0$ and then $y=0$.
\item [\textbf{ex.}]\textbf{${\displaystyle \lim_{(x,y)\rightarrow\left(0,0\right)}\frac{y^{3}+x^{3}}{xy^{2}}}$},
use lines $x=y$ and then $x=-y$.\vfill
\end{itemize}
\end{itemize}

\subsubsection*{13.4 Partial Derivatives}
\begin{itemize}
\item Given functions $f(x,y)$, find $f_{x}$, $f_{y}$, $f_{xx}$, $f_{yy}$,
$f_{xy}$, and $f_{yx}$. 

\begin{itemize}
\item [\textbf{ex.}]$f(x,y)=e^{x+7},\,\cos xy,\,2x^{3}+3y^{2}+1$\vfill
\end{itemize}
\end{itemize}

\subsubsection*{13.5 The Chain Rule}
\begin{itemize}
\item Use the Chain rule to find the derivatives of functions with one independent
variable (only two-variable functions). 

\begin{itemize}
\item [\textbf{ex.}]Find $dz/dt$ where $z=x\sin y$, $x=t^{2}$, and $y=4t^{3}$\vfill
\end{itemize}
\end{itemize}

\subsubsection*{13.6 Directional Derivatives and the Gradient}
\begin{itemize}
\item Compute the gradient of two-variable functions.

\begin{itemize}
\item [\textbf{ex.}]Let \textbf{$f(x,y)=4x^{2}-2xy+y^{2}$}. Find $\nabla f(x,y)$
at $P(-1,-5)$. Equivalently, find the steepest ascent of $f(x,y)$
at $P(-1,-5)$.\vfill
\end{itemize}
\item Find a unit vector that points in the direction of the steepest ascent
and descent of a function.
\item Find a vector that gives a direction of no change of a function.

\begin{itemize}
\item [\textbf{ex.}]For $f(x,y)=x^{2}-4y^{2}-9$ at $P(1,-2)$\vfill
\end{itemize}
\end{itemize}

\subsubsection*{13.8 Maximum and Minimum Problems}
\begin{itemize}
\item Find the critical points of a function, and use the Second Derivative
Test to determine (when possible) whether each critical point corresponds
to a local maximum, minimum, or saddle point.

\begin{itemize}
\item [\textbf{ex.}]$f(x,y)=-4x^{2}+8y^{2}-3$, $f(x,y)=\sqrt{x^{2}+y^{2}-4x+5}$,
$f(x,y)=ye^{x}-e^{y}$.\vfill
\end{itemize}
\end{itemize}

\end{document}
